\documentclass[letterpaper,12pt]{article}
\usepackage[margin=.5in]{geometry}
\usepackage{tikz}
\usepackage{helvet}
\renewcommand{\familydefault}{\sfdefault}
\pagestyle{empty}
\pdfmapfile{+cm.map}
\pdfmapfile{+ps2pk35.map}
\setlength{\parindent}{0pt}
\hyphenpenalty=10000
\exhyphenpenalty=10000
\begin{document}
\null\begin{tikzpicture}[remember picture,overlay,x=1cm,y=1cm]
\begin{scope}[shift={(current page.north west)}]
\node[anchor=north west,inner sep=0pt,font=\fontsize{11.5}{13}\selectfont] at (1.4,-.72) {Week 30 / Two-pan weight kits / Grades 4--5};
\node[anchor=south west,inner sep=0pt,font=\fontsize{9}{11}\selectfont] at (1.4,-27.2) {Bellingham Math Circle / Week 30 / F30-45-v2};
\node[anchor=south east,inner sep=0pt,font=\fontsize{9}{11}\selectfont] at (20.15,-27.2) {1};
\begin{scope}[shift={(1.4,-1.65)}]
\node[anchor=north west,inner sep=0pt,align=left,text width=18.5cm,font=\fontsize{13}{16.38}\selectfont] at (0,0) {Keep the target on the left. Each weight may go on either pan or stay off. Use each weight at most once. Equal totals balance.};
\node[anchor=north west,inner sep=0pt,align=left,text width=18.5cm,font=\fontsize{12}{15.120000000000001}\selectfont] at (0,-1.5) {Target 3; weights 1 and 4};
\draw[rounded corners=5pt,line width=.7pt] (0.0,-2.2) rectangle (2.65,-4.45);
\node[font=\fontsize{15}{17}\selectfont,inner sep=1pt] at (1.32,-2.62) {1};
\fill (0.65,-3.25) circle (1.5pt);
\draw[rounded corners=5pt,line width=.7pt] (3.3,-2.2) rectangle (5.949999999999999,-4.45);
\node[font=\fontsize{15}{17}\selectfont,inner sep=1pt] at (4.62,-2.62) {4};
\fill (3.9499999999999997,-3.25) circle (1.5pt);
\fill (4.6,-3.25) circle (1.5pt);
\fill (5.25,-3.25) circle (1.5pt);
\fill (3.9499999999999997,-3.63) circle (1.5pt);
\draw[->,line width=.8pt] (8.0,-3.2) -- (9.3,-3.2);
\node[font=\fontsize{12}{14}\selectfont,inner sep=1pt] at (12.4,-3.1) {put 1 beside the target};
\draw[rounded corners=9pt,line width=.7pt] (0,-5.1) rectangle (7.9,-8.2);
\draw[rounded corners=9pt,line width=.7pt] (9.1,-5.1) rectangle (17.0,-8.2);
\draw[fill=gray!12,line width=.7pt] (0.25,-5.3999999999999995) rectangle (2.45,-7.8999999999999995);
\node[font=\fontsize{10}{12}\selectfont,inner sep=1pt] at (1.35,-6.05) {target};
\node[font=\fontsize{16}{18}\selectfont,inner sep=1pt] at (1.35,-7.0) {3};
\draw (3.5,-6.6499999999999995) circle (0.8cm);
\node[font=\fontsize{14}{16}\selectfont,inner sep=1pt] at (3.5,-6.6499999999999995) {1};
\draw (10.65,-6.6499999999999995) circle (0.8cm);
\node[font=\fontsize{14}{16}\selectfont,inner sep=1pt] at (10.65,-6.6499999999999995) {4};
\node[font=\fontsize{18}{20}\selectfont,inner sep=1pt] at (8.5,-6.6499999999999995) {=};
\node[font=\fontsize{13}{15}\selectfont,inner sep=1pt] at (8.5,-8.8) {3 + 1 = 4};
\node[anchor=north west,inner sep=0pt,align=left,text width=18.6cm,font=\fontsize{14}{17.78}\selectfont] at (0,-10) {\textbf{Problem 1:} Choose two different whole-number weights that balance every target from 1 through 4. Find every possible kit.};
\draw[rounded corners=9pt,line width=.7pt] (0,-13.0) rectangle (7.9,-16.1);
\draw[rounded corners=9pt,line width=.7pt] (9.1,-13.0) rectangle (17.0,-16.1);
\draw[fill=gray!12,line width=.7pt] (0.25,-13.3) rectangle (2.45,-15.8);
\node[font=\fontsize{10}{12}\selectfont,inner sep=1pt] at (1.35,-13.95) {target};
\node[font=\fontsize{16}{18}\selectfont,inner sep=1pt] at (1.35,-14.9) {2};
\node[font=\fontsize{18}{20}\selectfont,inner sep=1pt] at (8.5,-14.55) {=};
\draw[rounded corners=9pt,line width=.7pt] (0,-17.6) rectangle (7.9,-20.700000000000003);
\draw[rounded corners=9pt,line width=.7pt] (9.1,-17.6) rectangle (17.0,-20.700000000000003);
\draw[fill=gray!12,line width=.7pt] (0.25,-17.900000000000002) rectangle (2.45,-20.400000000000002);
\node[font=\fontsize{10}{12}\selectfont,inner sep=1pt] at (1.35,-18.55) {target};
\node[font=\fontsize{16}{18}\selectfont,inner sep=1pt] at (1.35,-19.5) {4};
\node[font=\fontsize{18}{20}\selectfont,inner sep=1pt] at (8.5,-19.150000000000002) {=};
\draw[gray!45,line width=.35pt] (0,-23.0) -- (18.5,-23.0);\end{scope}\end{scope}\end{tikzpicture}
\newpage
\null\begin{tikzpicture}[remember picture,overlay,x=1cm,y=1cm]
\begin{scope}[shift={(current page.north west)}]
\node[anchor=north west,inner sep=0pt,font=\fontsize{11.5}{13}\selectfont] at (1.4,-.72) {Week 30 / Two-pan weight kits / Grades 4--5};
\node[anchor=south west,inner sep=0pt,font=\fontsize{9}{11}\selectfont] at (1.4,-27.2) {Bellingham Math Circle / Week 30 / F30-45-v2};
\node[anchor=south east,inner sep=0pt,font=\fontsize{9}{11}\selectfont] at (20.15,-27.2) {2};
\begin{scope}[shift={(1.4,-1.65)}]
\node[anchor=north west,inner sep=0pt,align=left,text width=18.6cm,font=\fontsize{14}{17.78}\selectfont] at (0,0) {\textbf{Problem 2:} Add one weight to 1 and 3. Which choice balances the longest unbroken run of targets starting at 1? Explain why a smaller or larger choice cannot do better.};
\draw[black,line width=.65pt] (0.0,-3.0) rectangle (1.8,-4.2);
\node[font=\fontsize{16}{18}\selectfont,inner sep=1pt] at (0.9,-3.6) {1};
\draw[black,line width=.65pt] (2.2,-3.0) rectangle (4.0,-4.2);
\node[font=\fontsize{16}{18}\selectfont,inner sep=1pt] at (3.1,-3.6) {2};
\draw[black,line width=.65pt] (4.4,-3.0) rectangle (6.2,-4.2);
\node[font=\fontsize{16}{18}\selectfont,inner sep=1pt] at (5.300000000000001,-3.6) {3};
\draw[black,line width=.65pt] (6.6000000000000005,-3.0) rectangle (8.4,-4.2);
\node[font=\fontsize{16}{18}\selectfont,inner sep=1pt] at (7.500000000000001,-3.6) {4};
\draw[black,line width=.65pt] (8.8,-3.0) rectangle (10.600000000000001,-4.2);
\node[font=\fontsize{16}{18}\selectfont,inner sep=1pt] at (9.700000000000001,-3.6) {5};
\draw[black,line width=.65pt] (11.0,-3.0) rectangle (12.8,-4.2);
\node[font=\fontsize{16}{18}\selectfont,inner sep=1pt] at (11.9,-3.6) {6};
\draw[black,line width=.65pt] (13.200000000000001,-3.0) rectangle (15.000000000000002,-4.2);
\node[font=\fontsize{16}{18}\selectfont,inner sep=1pt] at (14.100000000000001,-3.6) {7};
\draw[black,line width=.65pt] (15.400000000000002,-3.0) rectangle (17.200000000000003,-4.2);
\node[font=\fontsize{16}{18}\selectfont,inner sep=1pt] at (16.3,-3.6) {8};
\draw[black,line width=.65pt] (0.0,-4.6) rectangle (1.8,-5.8);
\node[font=\fontsize{16}{18}\selectfont,inner sep=1pt] at (0.9,-5.199999999999999) {9};
\draw[black,line width=.65pt] (2.2,-4.6) rectangle (4.0,-5.8);
\node[font=\fontsize{16}{18}\selectfont,inner sep=1pt] at (3.1,-5.199999999999999) {10};
\draw[black,line width=.65pt] (4.4,-4.6) rectangle (6.2,-5.8);
\node[font=\fontsize{16}{18}\selectfont,inner sep=1pt] at (5.300000000000001,-5.199999999999999) {11};
\draw[black,line width=.65pt] (6.6000000000000005,-4.6) rectangle (8.4,-5.8);
\node[font=\fontsize{16}{18}\selectfont,inner sep=1pt] at (7.500000000000001,-5.199999999999999) {12};
\draw[black,line width=.65pt] (8.8,-4.6) rectangle (10.600000000000001,-5.8);
\node[font=\fontsize{16}{18}\selectfont,inner sep=1pt] at (9.700000000000001,-5.199999999999999) {13};
\draw[black,line width=.65pt] (11.0,-4.6) rectangle (12.8,-5.8);
\node[font=\fontsize{16}{18}\selectfont,inner sep=1pt] at (11.9,-5.199999999999999) {14};
\draw[black,line width=.65pt] (13.200000000000001,-4.6) rectangle (15.000000000000002,-5.8);
\node[font=\fontsize{16}{18}\selectfont,inner sep=1pt] at (14.100000000000001,-5.199999999999999) {15};
\draw[rounded corners=9pt,line width=.7pt] (0,-8.0) rectangle (7.9,-11.1);
\draw[rounded corners=9pt,line width=.7pt] (9.1,-8.0) rectangle (17.0,-11.1);
\draw[fill=gray!12,line width=.7pt] (0.25,-8.3) rectangle (2.45,-10.8);
\node[font=\fontsize{10}{12}\selectfont,inner sep=1pt] at (1.35,-8.95) {target};
\node[font=\fontsize{16}{18}\selectfont,inner sep=1pt] at (1.35,-9.9) {5};
\node[font=\fontsize{18}{20}\selectfont,inner sep=1pt] at (8.5,-9.55) {=};
\draw[rounded corners=9pt,line width=.7pt] (0,-12.6) rectangle (7.9,-15.7);
\draw[rounded corners=9pt,line width=.7pt] (9.1,-12.6) rectangle (17.0,-15.7);
\draw[fill=gray!12,line width=.7pt] (0.25,-12.9) rectangle (2.45,-15.4);
\node[font=\fontsize{10}{12}\selectfont,inner sep=1pt] at (1.35,-13.549999999999999) {target};
\node[font=\fontsize{16}{18}\selectfont,inner sep=1pt] at (1.35,-14.5) {7};
\node[font=\fontsize{18}{20}\selectfont,inner sep=1pt] at (8.5,-14.15) {=};
\draw[gray!45,line width=.35pt] (0,-20.0) -- (18.5,-20.0);
\draw[gray!45,line width=.35pt] (0,-21.15) -- (18.5,-21.15);
\draw[gray!45,line width=.35pt] (0,-22.3) -- (18.5,-22.3);\end{scope}\end{scope}\end{tikzpicture}
\newpage
\null\begin{tikzpicture}[remember picture,overlay,x=1cm,y=1cm]
\begin{scope}[shift={(current page.north west)}]
\node[anchor=north west,inner sep=0pt,font=\fontsize{11.5}{13}\selectfont] at (1.4,-.72) {Week 30 / Two-pan weight kits / Grades 4--5};
\node[anchor=south west,inner sep=0pt,font=\fontsize{9}{11}\selectfont] at (1.4,-27.2) {Bellingham Math Circle / Week 30 / F30-45-v2};
\node[anchor=south east,inner sep=0pt,font=\fontsize{9}{11}\selectfont] at (20.15,-27.2) {3};
\begin{scope}[shift={(1.4,-1.65)}]
\node[anchor=north west,inner sep=0pt,align=left,text width=18.6cm,font=\fontsize{14}{17.78}\selectfont] at (0,0) {\textbf{Problem 3:} Can any positive target be balanced in two different ways with weights 1, 3 and 9? Decide for every target this kit can balance.};
\draw[rounded corners=9pt,line width=.7pt] (0,-3.0) rectangle (7.9,-6.1);
\draw[rounded corners=9pt,line width=.7pt] (9.1,-3.0) rectangle (17.0,-6.1);
\draw[fill=gray!12,line width=.7pt] (0.25,-3.3) rectangle (2.45,-5.8);
\node[font=\fontsize{10}{12}\selectfont,inner sep=1pt] at (1.35,-3.95) {target};
\node[font=\fontsize{16}{18}\selectfont,inner sep=1pt] at (1.35,-4.9) {4};
\node[font=\fontsize{18}{20}\selectfont,inner sep=1pt] at (8.5,-4.55) {=};
\draw[rounded corners=9pt,line width=.7pt] (0,-7.6) rectangle (7.9,-10.7);
\draw[rounded corners=9pt,line width=.7pt] (9.1,-7.6) rectangle (17.0,-10.7);
\draw[fill=gray!12,line width=.7pt] (0.25,-7.8999999999999995) rectangle (2.45,-10.399999999999999);
\node[font=\fontsize{10}{12}\selectfont,inner sep=1pt] at (1.35,-8.549999999999999) {target};
\node[font=\fontsize{16}{18}\selectfont,inner sep=1pt] at (1.35,-9.5) {7};
\node[font=\fontsize{18}{20}\selectfont,inner sep=1pt] at (8.5,-9.15) {=};
\draw[gray!45,line width=.35pt] (0,-14.0) -- (18.5,-14.0);
\draw[gray!45,line width=.35pt] (0,-15.15) -- (18.5,-15.15);
\draw[gray!45,line width=.35pt] (0,-16.3) -- (18.5,-16.3);
\draw[gray!45,line width=.35pt] (0,-17.45) -- (18.5,-17.45);
\draw[gray!45,line width=.35pt] (0,-18.6) -- (18.5,-18.6);
\draw[gray!45,line width=.35pt] (0,-19.75) -- (18.5,-19.75);
\draw[gray!45,line width=.35pt] (0,-20.9) -- (18.5,-20.9);\end{scope}\end{scope}\end{tikzpicture}
\newpage
\null\begin{tikzpicture}[remember picture,overlay,x=1cm,y=1cm]
\begin{scope}[shift={(current page.north west)}]
\node[anchor=north west,inner sep=0pt,font=\fontsize{11.5}{13}\selectfont] at (1.4,-.72) {Week 30 / Two-pan weight kits / Grades 4--5};
\node[anchor=south west,inner sep=0pt,font=\fontsize{9}{11}\selectfont] at (1.4,-27.2) {Bellingham Math Circle / Week 30 / F30-45-v2};
\node[anchor=south east,inner sep=0pt,font=\fontsize{9}{11}\selectfont] at (20.15,-27.2) {4};
\begin{scope}[shift={(1.4,-1.65)}]
\node[anchor=north west,inner sep=0pt,align=left,text width=18.6cm,font=\fontsize{14}{17.78}\selectfont] at (0,0) {\textbf{Problem 4:} Could any three-weight kit balance every target from 1 through 14? Find one or explain why none can work.};
\draw[gray!45,line width=.35pt] (0,-3.0) -- (18.5,-3.0);
\draw[gray!45,line width=.35pt] (0,-4.15) -- (18.5,-4.15);
\draw[gray!45,line width=.35pt] (0,-5.3) -- (18.5,-5.3);
\draw[gray!45,line width=.35pt] (0,-6.449999999999999) -- (18.5,-6.449999999999999);
\draw[gray!45,line width=.35pt] (0,-7.6) -- (18.5,-7.6);
\draw[gray!45,line width=.35pt] (0,-8.75) -- (18.5,-8.75);
\draw[gray!45,line width=.35pt] (0,-9.899999999999999) -- (18.5,-9.899999999999999);
\draw[gray!45,line width=.35pt] (0,-11.049999999999999) -- (18.5,-11.049999999999999);
\node[anchor=north west,inner sep=0pt,align=left,text width=18.6cm,font=\fontsize{14}{17.78}\selectfont] at (0,-15) {\textbf{Problem 5:} Choose three different whole-number weights. How many different positive targets can your kit balance? Can another kit balance more, even if there are gaps?};
\draw[gray!45,line width=.35pt] (0,-18.0) -- (18.5,-18.0);
\draw[gray!45,line width=.35pt] (0,-19.15) -- (18.5,-19.15);
\draw[gray!45,line width=.35pt] (0,-20.3) -- (18.5,-20.3);
\draw[gray!45,line width=.35pt] (0,-21.45) -- (18.5,-21.45);
\draw[gray!45,line width=.35pt] (0,-22.6) -- (18.5,-22.6);\end{scope}\end{scope}\end{tikzpicture}
\newpage
\null\begin{tikzpicture}[remember picture,overlay,x=1cm,y=1cm]
\begin{scope}[shift={(current page.north west)}]
\node[anchor=north west,inner sep=0pt,font=\fontsize{11.5}{13}\selectfont] at (1.4,-.72) {Week 30 / Two-pan weight kits / Grades 4--5};
\node[anchor=south west,inner sep=0pt,font=\fontsize{9}{11}\selectfont] at (1.4,-27.2) {Bellingham Math Circle / Week 30 / F30-45-v2};
\node[anchor=south east,inner sep=0pt,font=\fontsize{9}{11}\selectfont] at (20.15,-27.2) {5};
\begin{scope}[shift={(1.4,-1.65)}]
\node[anchor=north west,inner sep=0pt,align=left,text width=18.6cm,font=\fontsize{14}{17.78}\selectfont] at (0,0) {\textbf{Problem 6:} Use the kit 1, 3 and 9. Add one weight so that the new kit balances every target in the longest possible run starting at 1, each in exactly one way. Explain why there are neither gaps nor repeated answers.};
\draw[gray!45,line width=.35pt] (0,-3.5) -- (18.5,-3.5);
\draw[gray!45,line width=.35pt] (0,-4.65) -- (18.5,-4.65);
\draw[gray!45,line width=.35pt] (0,-5.8) -- (18.5,-5.8);
\draw[gray!45,line width=.35pt] (0,-6.949999999999999) -- (18.5,-6.949999999999999);
\draw[gray!45,line width=.35pt] (0,-8.1) -- (18.5,-8.1);
\draw[gray!45,line width=.35pt] (0,-9.25) -- (18.5,-9.25);
\draw[gray!45,line width=.35pt] (0,-10.399999999999999) -- (18.5,-10.399999999999999);
\node[anchor=north west,inner sep=0pt,align=left,text width=18.6cm,font=\fontsize{14}{17.78}\selectfont] at (0,-14) {\textbf{Problem 7:} With five weights, what is the longest run of whole-number targets starting at 1 that any kit could balance? Give a kit that reaches your bound, and explain why it works.};
\draw[gray!45,line width=.35pt] (0,-17.0) -- (18.5,-17.0);
\draw[gray!45,line width=.35pt] (0,-18.15) -- (18.5,-18.15);
\draw[gray!45,line width=.35pt] (0,-19.3) -- (18.5,-19.3);
\draw[gray!45,line width=.35pt] (0,-20.45) -- (18.5,-20.45);
\draw[gray!45,line width=.35pt] (0,-21.6) -- (18.5,-21.6);
\draw[gray!45,line width=.35pt] (0,-22.75) -- (18.5,-22.75);\end{scope}\end{scope}\end{tikzpicture}
\end{document}
